\documentclass[11pt]{article}
\usepackage[margin=1in]{geometry}
\usepackage{hyperref}
\usepackage{booktabs}

\title{FREUID Challenge 2026 Technical Report}
\author{Hanif Noer Rofiq}
\date{July 2026}

\begin{document}
\maketitle

\section{Introduction}
This report describes the final FREUID Challenge 2026 submission. The task is binary fraud detection for identity-document images. Scores are numeric fraud probabilities, where higher values indicate higher confidence that a document is fraudulent.

\section{Method}
The final model is based on \texttt{convnextv2\_large.fcmae\_ft\_in22k\_in1k\_384} from \texttt{timm}, fine-tuned as a single-logit binary classifier.

Images are resized/cropped to \texttt{672x672}. Training augmentations used random resized crop with scale \texttt{0.88--1.0}, color jitter, light Gaussian blur, and small rotation. Inference uses resize to \texttt{672}, center crop, ImageNet normalization, and sigmoid conversion from logits to fraud probabilities.

\section{Data}
The model uses the official FREUID Challenge training data. The local training labels contain 69,352 rows, with 40,005 negative and 29,347 positive labels.

External resources:
\begin{itemize}
  \item \texttt{timm} pretrained ConvNeXtV2 weights.
  \item PyTorch, torchvision, Pillow, and timm runtime dependencies.
  \item Final checkpoint hosted in a public Google Drive folder for Docker build-time retrieval.
\end{itemize}

\section{Training}
Training was performed in Google Colab/GPU environment using the notebooks in this repository.

\begin{itemize}
  \item Main notebook: \texttt{notebooks/freuid\_convnextv2\_g4\_672\_alltrain\_best\_public.ipynb}
  \item Batch size: 24
  \item Initial learning rate: \texttt{1e-5}
  \item Weight decay: \texttt{0.05}
  \item Training epochs: 3
  \item Final selected checkpoint: epoch 3 \texttt{best.pt} from the run that produced the public score below
\end{itemize}

\section{Results}
The current selected model is the user-reported best public leaderboard submission.

\begin{center}
\begin{tabular}{ll}
\toprule
Item & Value \\
\midrule
Final submission label/time & \texttt{submission\_convnextv2\_g4\_672\_alltrain.csv} / TODO time \\
Public score & 0.00267 \\
Private score & TODO \\
Selected checkpoint SHA256 & TODO \\
\bottomrule
\end{tabular}
\end{center}

\section{Inference}
The Docker package reads a flat image directory mounted at \texttt{/data} and writes \texttt{/submissions/submission.csv}. The row id is the filename stem. The Docker runtime does not require network access.

\section{Reproducibility}
Build and run. The Docker build downloads the final checkpoint from the
documented public Google Drive folder if it is not already present under
\texttt{reproducibility/model/model.pt}. Inference itself uses no network.

\begin{verbatim}
cd reproducibility
docker build -t freuid-repro:local .
mkdir -p out
docker run --rm --network none \
  -v /path/to/flat/test/images:/data:ro \
  -v "$PWD/out:/submissions" \
  freuid-repro:local
\end{verbatim}

If the model checkpoint is not baked into the image, mount it read-only:

\begin{verbatim}
docker run --rm --network none \
  -v /path/to/flat/test/images:/data:ro \
  -v /path/to/final_checkpoint.pt:/models/model.pt:ro \
  -v "$PWD/out:/submissions" \
  freuid-repro:local
\end{verbatim}

\end{document}
